\clearpage
\chapter*{ABSTRAK}
\addcontentsline{toc}{chapter}{ABSTRAK}
\begin{center}
	\center
	\begin{singlespace}
		\large\bfseries\MakeUppercase{\thetitle}

		\normalfont\normalsize
		Oleh:

		\bfseries \theauthor
	\end{singlespace}
\end{center}

\begin{singlespace}
	\small
	Transformasi Rekam Medis Elektronik (RME) menuntut mekanisme kontrol akses yang menjaga kendali pasien sekaligus mendukung alur pelayanan kesehatan kolaboratif. DecMed sebagai sistem RME terdesentralisasi telah memanfaatkan IOTA, IPFS, dan Proxy Re-Encryption (PRE), tetapi mekanisme sebelumnya masih bergantung pada pasien untuk menerbitkan akses baru, masih bersifat \textit{coarse-grained} karena hanya berbasis \textit{role} dan \textit{purpose}, serta menyimpan data RME sebagai objek besar yang sulit ditegakkan sesuai kebutuhan layanan. Kontrol akses pada DecMed diperluas agar keputusan akses mencakup aktor, objek RME, hak akses, dan kondisi akses secara lebih luas, sekaligus mendukung segmentasi data, delegasi terbatas, revokasi, audit, dan verifikasi identitas aktor pada rantai delegasi. Pendekatan yang digunakan mencakup penerapan \textit{Policy-Based Access Control} (PBAC) sebagai model kebijakan, Macaroon sebagai media pembawa kapabilitas dengan \textit{caveats} untuk atenuasi, segmentasi data RME \textit{off-chain} berdasarkan kategori dataset dan fungsi klinis, pencatatan metadata serta audit pada IOTA, dan PRE sebagai titik penegakan kebijakan sebelum data terenkripsi diambil atau diproses ulang untuk penerima akses. Implementasi mencakup komponen segmentasi RME, token Macaroon, penegakan otoritas pada \textit{Server} PRE, \textit{smart contract} IOTA Move, \textit{Patient Client}, dan \textit{Hospital Client} untuk alur pasien, petugas rekam medis, perawat, dokter, petugas laboratorium, dan apoteker. Pengujian dilakukan melalui skenario fungsional \texttt{P-KF-01} sampai \texttt{P-KF-12} pada alur rawat jalan dan rawat inap serta pengujian nonfungsional \texttt{P-KNF-01} sampai \texttt{P-KNF-13}. Evaluasi menunjukkan seluruh kebutuhan fungsional dan nonfungsional yang dirumuskan dinyatakan terpenuhi. Hasil penelitian ini menunjukkan bahwa kontrol akses berbasis kebijakan, segmentasi RME, delegasi terbatas, revokasi, audit, dan verifikasi identitas aktor pada rantai delegasi dapat dioperasionalkan dalam DecMed pada batas implementasi prototipe yang diuji.

	\textbf{\textit{Kata kunci: Rekam Medis Elektronik, Policy-Based Access Control, Macaroon, IOTA, Proxy Re-Encryption}}

\end{singlespace}
\clearpage